\subsection{Public Evidence and Reconstruction}
\label{sec:audit}
E5 illustrates why a candidate match can leave its history unresolved.
The six-release audit asks the complementary question of whether public
artifacts support reconstruction of a published agent outcome. It applies a
predeclared nine-field rubric to six public releases, yielding
$6\times9=54$ release--field assessments (A tier); the supplement retains
the source-linked method, dispositions, and mechanism examples.

No release satisfies all nine fields (0/6): the retained public material does
not bind implementation, configuration, task, immutable environment, execution
identity, and result for one end-to-end success. The 54 field-level cells carry graded rather than uniformly absent
evidence, with the per-cell dispositions in the supplement.
The published-outcome field alone blocks the all-nine conjunction;
Table~\ref{tab:audit-fields} therefore reports each field separately.
Removing it still leaves 0/6 satisfying all eight remaining fields; this
sensitivity check does not validate the rubric or make every field
necessary for every claim. Clock and cost evidence support timing and resource
claims; weakness and discovery fields support mechanism and autonomy claims.
They are not prerequisites for a candidate-correctness claim. The all-nine
conjunction addresses complete historical reconstruction only, and the finding
concerns reconstructability of the published evidence rather than correctness
of the audited systems. The audit remains separate from V4, E5, and the
executed attack families.

\begin{table}[t]
\centering\footnotesize
\caption{Release-audit field profiles at the inspected snapshots. Y/P/N are
yes/partial/no under the frozen evidence rules, not quality scores.}
\label{tab:audit-fields}
\begin{tabularx}{\columnwidth}{@{}Y c c c@{}}
\toprule
\tablehead Field and supported claim & Y & P & N\\
\midrule
Clock boundary / elapsed time & 1 & 3 & 2\\
Oracle isolation / grading authority & 0 & 4 & 2\\
Failure retention / outcome denominator & 0 & 1 & 5\\
Ordered traces / execution sequence & 0 & 1 & 5\\
Hash pinning / artifact identity & 0 & 1 & 5\\
Weakness attribution / mechanism & 0 & 2 & 4\\
Discovery--execution split / autonomy & 5 & 1 & 0\\
Query accounting / resource cost & 2 & 4 & 0\\
Published-outcome probe / reconstruction & 0 & 1 & 5\\
\bottomrule
\end{tabularx}
\end{table}

\subsection{Attribution Policy and Tests}\label{sec:attribution-contract}
The audit concerns available records; RQ3 additionally requires testing what
records justify. The policy separates authenticity from consistency: a signed
record must come from an authorized role, and its contents must agree with
the other records for the requested claim. An observer-signed record
covers the exact executed bytes and captured stdout; a reference-signed record
recomputes the task predicate from private target state against that stdout and
names the observer record. Attribution requires agreement of the declared
digests, matching identity fields, authorized observer and reference roles, and
the applicable replay policy; the claimant's own \texttt{claims} block never
influences the decision. Trust assumptions concern the observer, reference
verifier, and role policy.

Whether a non-zero exit may accompany a predicate-satisfying output is a
declared policy choice, and two readings were fixed before evaluation. Across
four corpora, the property check evaluates 1{,}407 mutants under each reading,
for 11{,}256 evaluations without a soundness violation: bounded mechanized
testing, not a formal proof. The live-target and
loopback-ingress exercises test the same rule against fabricated claims; the
supplement gives their controls and outcomes.

The contract uses existing provenance mechanisms to express the policy: a
configured in-toto inspection reproduces 21/22 decisions, reaching 22/22 with a
shared replay ledger, while the SLSA arm stops at provenance acquisition and is
not comparable. These results support an explicit claim policy and clarify its
configuration requirements; they do not establish superiority over provenance
frameworks.
